
% Dateiname: MoE_Glossar_Pfad_261029-2.tex

\documentclass[11pt,a4paper]{article}

\usepackage[T1]{fontenc}
\usepackage[utf8]{inputenc}
\usepackage[ngerman]{babel}
\usepackage[a4paper,margin=25mm]{geometry}
\usepackage{amsmath,amssymb}
\usepackage{helvet}
\usepackage{sfmath}
\renewcommand{\familydefault}{\sfdefault}
\usepackage{microtype}
\usepackage[hidelinks]{hyperref}
\usepackage{lastpage}
\usepackage{fancyhdr}
\usepackage{parskip}

\pagestyle{fancy}
\fancyhf{}
\fancyfoot[C]{\thepage\ (\pageref{LastPage})}
\renewcommand{\headrulewidth}{0pt}
\renewcommand{\footrulewidth}{0pt}

\hypersetup{
  pdftitle={MoE-Glossar: Pfad},
  pdfauthor={}
}

\begin{document}

\vspace*{-1.8cm}

\noindent
\begin{minipage}[t]{0.66\textwidth}
{\LARGE\bfseries Pfad /
Elementare Pfadtopologien}
\end{minipage}
\hfill
\begin{minipage}[t]{0.30\textwidth}
\raggedleft
{\small MoE-Glossareintrag}\\
{\small Version 261029-2}
\end{minipage}

\vspace{1.4em}

\section*{Definition}

Ein Pfad ist eine stabile,
geschlossene Apeiron-Strömungsform
innerhalb der räumlichen
Kontaktordnung der Dichten
QUBLE-Packung (DQP).

Die räumliche Führung der
Strömung bestimmt die Geometrie
und Topologie des Pfades.

\section*{Elementare Pfadformen}

Im MoE werden drei elementare
Pfadformen beschrieben:

\begin{itemize}
\item Ringpfad,
\item Tennisball-Pfad (TBTP),
\item Kleeblattknoten
      (engl.\ Trefoil Knot).
\end{itemize}

Die Pfadformen sind den
elementaren Teilchenstrukturen
zugeordnet.

\section*{Ringpfad}

Der Ringpfad ist die
Pfadstruktur des Photons
$\gamma$.

Er bildet eine geschlossene
ringförmige Strömungsbahn.

Sein Pfadlängenfaktor ist:

\[
L_P(\gamma)=1.
\]

\section*{Tennisball-Pfad}

Der Tennisball-Pfad (TBTP)
ist die Pfadstruktur von
Elektron und Positron sowie
Myon und Antimyon.

Er besteht aus vier
Halbkreisbögen auf einer
kugelförmigen Trägergeometrie.

Sein Pfadlängenfaktor ist:

\[
L_P(\mathrm{TBTP})=2.
\]

Die Pfadführung von Elektron
und Positron besitzt eine
Elektronachse.

Bei Myon und Antimyon
erfolgt ein Achsenwechsel.
Die Myonachse berührt die
Tennisballkurve.

\section*{Kleeblattknoten}

Der Kleeblattknoten
(Trefoil Knot) bildet die
Pfadstruktur von Neutrino
und Antineutrino.

Er besitzt eine geschlossene,
dreifach verschlungene
Pfadtopologie.

Sein Pfadlängenfaktor ist:

\[
L_P(\nu)=3.
\]

\section*{Pfadlänge}

Die Pfadlänge $s_P$ ist die
geometrische Länge des
vollständigen geschlossenen
Pfadverlaufs.

Sie ergibt sich aus der
räumlichen Geometrie
des jeweiligen Pfades.

Der Pfadlängenfaktor $L_P$
bezeichnet die normierte
Pfadlänge der elementaren
Pfadformen.

Der Ringpfad des Photons
dient als Bezugsgröße:

\[
L_P(\gamma)=1.
\]

Für einen kreisförmigen
Ringpfad mit dem Durchmesser
$D_\gamma$ gilt:

\[
s_\gamma=\pi D_\gamma.
\]

Mit der im MoE festgelegten
Beziehung

\[
D_\gamma=\frac{\lambda}{\pi}
\]

ergibt sich:

\[
s_\gamma=\lambda.
\]

Die Pfadlängenfaktoren
der elementaren Pfade
lauten:

\[
\begin{aligned}
L_P(\gamma) &= 1,\\
L_P(e^\pm) &= 2,\\
L_P(\mu^\pm) &= 2,\\
L_P(\nu) &= 3.
\end{aligned}
\]

\section*{Pfadlänge und QUBLE-Zahl}

Die QUBLE-Zahl $N_Q$
bezeichnet die Anzahl
der einem Pfad zugeordneten
QUBLES.

Für die elementaren
Pfadstrukturen gelten
die Beziehungen:

\[
N_Q(e)=2N_Q(\gamma),
\]

\[
N_Q(\nu)=3N_Q(\gamma).
\]

Die Zuordnungen entsprechen
den festgelegten
Pfadlängenfaktoren.

\section*{Pfad und Packung}

Die DQP bildet die räumliche
Kontaktordnung, in der die
Pfadstrukturen geführt werden.

Die Geometrie der QUBLES,
ihrer Kontaktlagen und
Interstitien bestimmt die
räumlichen Bedingungen
der Pfadführung.

\section*{Emergenzfolge}

\[
\begin{gathered}
\text{Apeiron}
\rightarrow
\text{Expansion}
\rightarrow
\text{Quantelung}
\\[0.4em]
\rightarrow
\text{QUBLE}
\rightarrow
\text{SDQP}
\rightarrow
\text{DQP}
\rightarrow
\text{Pfad}
\end{gathered}
\]

\vspace{0.9em}
\hrule
\vspace{0.45em}

{\footnotesize
MoE-Glossar:
\emph{Apeiron und Apeiron-Expansion},
\emph{QUBLE},
\emph{SDQP},
\emph{DQP},
\emph{Interstitium},
\emph{Photon},
\emph{Elektron},
\emph{Myon},
\emph{Neutrino}.
}

\end{document}
